
\subsection{滤过 $A_\infty$ 同伦的构造 II：同伦的同伦之代数框架}
\label{subsec:7-2-12}

本小节及下一小节要使用大量记号。
不过，它们是前面各小节对一般
$A_\infty$ 或 $A_{n,K}$ 结构研究的自然延续；
为了在一般系数环 $R$ 上证明定理~\ref{thm:4-6-25}，
这些记号可能难以避免。

上一小节关于系数 $\mathbb Z_2$ 或 $\mathbb C$ 的结果，
已足以覆盖本书写作时作者所知的全部应用。
因此，主要关注弗洛尔理论应用的读者
可略过本节余部，直接读第 7.3 节。
第 7.2.14 小节仅在第 7.4 节使用。

本小节及下一小节将系统研究
\emph{同伦的同伦}。
它们用于从给定的滤过 $A_{n,K}$ 同伦，
在一般系数环 $R$ 上构造滤过 $A_\infty$ 同伦。
以下固定幺半群 $\mathfrak G$，
并假定所有滤过 $A_\infty$ 同态等
均为 $\mathfrak G$-有隙。
“滤过 $A_{n,K}$ 结构”等简写，
均指 $\mathfrak G$-滤过的相应结构。

本小节先建立所需的代数框架。
回忆前一小节的滤过 $A_\infty$ 同态构造
用到序列 $\bar n(\lambda_i)$ 与
$\bar k(\lambda_i)$，并要求
\[
\begin{array}{lll}
\lambda_i\to\infty,&
\bar n(\lambda_i)\to\infty,&
\bar k(\lambda_i)\to\infty,\\
\lambda'_i\to\infty,&
\bar n(\lambda'_i)\to\infty,&
\bar k(\lambda'_i)\to\infty.
\end{array}
\]
另假定
\[
\bar n(\lambda_i)\le\bar n(\lambda'_i),
\qquad
\bar k(\lambda_i)\le\bar k(\lambda'_i).
\]

\begin{definition}{7.2.182}\label{def:7-2-182}
\noindent\textup{(1)} 给定 $\lambda_i\to\infty$，
并给定上述序列
$(\bar n,\bar k)=
(\bar n(\lambda_i),\bar k(\lambda_i))$。

\noindent\textup{(7.2.182.1)}
设 $C_i$ 是一列滤过
$A_{\bar n(\lambda_i),\bar k(\lambda_i)}$
代数。称一列
$A_{\bar n(\lambda_i),\bar k(\lambda_i)}$
同伦等价
$\{I_i\colon C_i\to C_{i+1}\}$
为 $A_{(\bar n,\bar k)}$ 代数的
\emph{等价有向系统}\liftIndex{I266}。
下文为简便，简称为
$A_{(\bar n,\bar k)}$ 代数的\emph{有向系统}；
本书不使用 $I_i$ 不是同伦等价的情形。

\noindent\textup{(7.2.182.2)}
给定两个 $A_{(\bar n,\bar k)}$ 代数的
有向系统 $\{I_i\colon C_i\to C_{i+1}\}$、
$\{I'_i\colon C'_i\to C'_{i+1}\}$。
从前者到后者的\emph{态射}\liftIndex{I267}
$\{f_i\colon C_i\to C'_i\}$，
定义为一列滤过
$A_{\bar n(\lambda_i),\bar k(\lambda_i)}$
同态 $f_i\colon C_i\to C'_i$，使
\begin{equation}
I'_i\circ f_i\sim f_{i+1}\circ I_i,
\qquad
\bigl(\sim\text{ 表示滤过 }
A_{\bar n(\lambda_i),\bar k(\lambda_i)}
\text{ 同伦}\bigr).
\tag{7.2.183}\label{eq:7-2-183}
\end{equation}
见图式~\ref{diag:7-2-3}。

\par\medskip\noindent\begin{minipage}{\textwidth}
\centering
\[
\begin{array}{ccccccccc}
C_1&\xrightarrow{I_1}&C_2&\xrightarrow{I_2}&\cdots
&\xrightarrow{I_{k-1}}&C_k&\xrightarrow{I_k}&\cdots\\[1ex]
\downarrow f_1&&\downarrow f_2&&&&\downarrow f_k&&\\[-0.3ex]
C'_1&\xrightarrow{I'_1}&C'_2&\xrightarrow{I'_2}&\cdots
&\xrightarrow{I'_{k-1}}&C'_k&\xrightarrow{I'_k}&\cdots
\end{array}
\]
\par{\renewcommand{\thefigure}{7.2.3}\refstepcounter{figure}
\normalfont\bfseries 图式~\thefigure\label{diag:7-2-3}}
\end{minipage}\par\medskip

\noindent\textup{(2)}
设 $(C,\{\mathfrak m_{k,\beta}\}_
{(\beta,k)\precsim(n,K)})$
为 $\mathfrak G$-有隙的滤过 $A_{n,K}$ 代数。
它向滤过 $A_\infty$ 代数的\emph{提升}\liftIndex{I268}
是 $(C,\{\widetilde{\mathfrak m}_{k,\beta}\})$，
满足在 $(\beta,k)\precsim(n,K)$ 时
$\widetilde{\mathfrak m}_{k,\beta}
=\mathfrak m_{k,\beta}$。
若 $(n,K)\precsim(n',K')$，
向滤过 $A_{n',K'}$ 代数的提升也同样定义。

\noindent\textup{(3)}
设 $\{\mathfrak f_{k,\beta}\}_
{(\beta,k)\precsim(n,K)}\colon C\to C'$
是两个滤过 $A_\infty$ 代数间的
$\mathfrak G$-有隙滤过 $A_{n,K}$ 同态。
其向滤过 $A_\infty$ 同态的提升，
是滤过 $A_\infty$ 同态
$\{\widetilde{\mathfrak f}_{k,\beta}\}$，
满足在 $(\beta,k)\precsim(n,K)$ 时
$\widetilde{\mathfrak f}_{k,\beta}
=\mathfrak f_{k,\beta}$。
向滤过 $A_{n',K'}$ 同态的提升也同样定义。

\noindent\textup{(4)}
滤过 $A_{n,K}$ 同伦向滤过
$A_\infty$ 同伦的提升也同样定义。
\end{definition}

第~\ref{subsec:7-2-8} 与
第~\ref{subsec:7-2-9} 小节末的论证蕴含下述结果。

\begin{lemma}{7.2.184}\label{lem:7-2-184}
设 $\{I_i\colon C_i\to C_{i+1}\}$、
$\{I'_i\colon C'_i\to C'_{i+1}\}$
为 $A_{(\bar n,\bar k)}$ 代数的有向系统，
$\{f_i\colon C_i\to C'_i\}$
为二者间的态射。则：

\noindent\textup{(7.2.185.1)}
可将 $C_i$ 提升为滤过 $A_\infty$ 代数，
将 $I_i$ 提升为它们之间的滤过
$A_\infty$ 同态；仍沿用原记号。
这样的 $A_\infty$ 代数系统及同伦等价，
称为有向系统 $\{I_i\colon C_i\to C_{i+1}\}$
的一个\emph{提升}。
对 $\{I'_i\colon C'_i\to C'_{i+1}\}$
也同样成立。

\noindent\textup{(7.2.185.2)}
对 (7.2.185.1) 中任意这样的提升，
每个 $f_i\colon C_i\to C'_i$
均可提升为滤过 $A_\infty$ 同态，
使图式~\ref{diag:7-2-3}
在滤过 $A_\infty$ 同伦意义下交换。
将提升后的同态记作 $\widetilde f_i$；
称系统 $\{\widetilde f_i\}$
为态射 $\{f_i\}$ 的一个\emph{提升}。
无混淆时，下文也常用 $f_i$ 代替
$\widetilde f_i$。
\end{lemma}

引理的证明使用与第~\ref{subsec:7-2-8}
及第~\ref{subsec:7-2-9} 小节相似的归纳论证，
另见引理 7.2.208。
由图式~\ref{diag:7-2-4} 的归纳构造，
还可证明 (7.2.185.1) 中系统的提升
在滤过 $A_\infty$ 同伦等价意义下唯一。

\par\medskip\noindent\begin{minipage}{\textwidth}
\centering
\[
\begin{array}{ccccccccc}
C_1&\xrightarrow{I_1}&C_2&\xrightarrow{I_2}&\cdots
&\xrightarrow{I_{k-1}}&C_k&\xrightarrow{I_k}&\cdots\\[1ex]
\downarrow i_1&&\downarrow i_2&&&&\downarrow i_k&&\\[-0.3ex]
C_1&\xrightarrow{I_1}&C_2&\xrightarrow{I_2}&\cdots
&\xrightarrow{I_{k-1}}&C_k&\xrightarrow{I_k}&\cdots
\end{array}
\]
\par{\renewcommand{\thefigure}{7.2.4}\refstepcounter{figure}
\normalfont\bfseries 图式~\thefigure\label{diag:7-2-4}}
\end{minipage}\par\medskip

具体说，设有两种按 (7.2.185.1)
提升这些结构的方法，
分别置于图式~\ref{diag:7-2-4} 的上下两行。
从作为滤过
$A_{\bar n(\lambda_1),\bar k(\lambda_1)}$
同态的 $i_1=\mathrm{id}$ 出发，
利用引理~\ref{lem:7-2-129} 归纳构造
滤过 $A_{\bar n(\lambda_i),\bar k(\lambda_i)}$
同态 $i_j$（$j\le i$），
使图式~\ref{diag:7-2-4} 前 $i$ 列
在该级同伦意义下交换。
以下固定 (7.2.185.1) 所构造的一个提升。

现在关心 (7.2.185.2) 中滤过
$A_\infty$ 同态 $\widetilde f_i$
是否在滤过 $A_\infty$ 同伦意义下唯一。
若直接模仿上述 (7.2.185.1) 的唯一性证明，
便会遇到困难；主要原因是
图式~\ref{diag:7-2-3} 仅同伦交换。
另一方面，对固定的 $n,K$，
$\widetilde f_i$ 的滤过 $A_{n,K}$ 同伦型
显然与提升 $\{f_i\}$ 的选择无关。
这又是上一小节末遇到的问题。
为处理这一点，我们将定义~\ref{def:7-2-182}
细化如下。

\begin{definition}{7.2.182bis}\label{def:7-2-182bis}
设 $\{f_i\colon C_i\to C'_i\}$
是从有向系统
$\{I_i\colon C_i\to C_{i+1}\}$
到 $\{I'_i\colon C'_i\to C'_{i+1}\}$
的一个态射。
设 $\mathfrak C_i^{(r)}$ 是
$[0,1]\times C'_i$ 的一个滤过
$A_{\bar n(\lambda_i),\bar k(\lambda_i)}$
代数模型。考虑一列滤过同阶同态
$\mathcal F_i\colon C_i\to
\mathfrak C_{i+1}^{(r)}$，满足
\begin{equation}
\operatorname{Eval}_{r=0}\circ\mathcal F_i
=I'_i\circ f_i,
\qquad
\operatorname{Eval}_{r=1}\circ\mathcal F_i
=f_{i+1}\circ I_i.
\tag{7.2.186}\label{eq:7-2-186}
\end{equation}
称 $\{\mathcal F_i\colon C_i\to
\mathfrak C_{i+1}^{(r)}\}$
为态射 $\{f_i\}$ 的
\emph{交换误差同伦}\liftIndex{I269}。
这里取 $r$ 而非 $s$ 作区间参数，
是为与下文记号保持一致。
这种 $\mathcal F_i$ 的存在正是条件 (7.2.183)；
这里的要点是具体选定这些同伦。
\end{definition}

于是可将引理~\ref{lem:7-2-184} 推广如下。

\begin{statement}{引理}{7.2.184bis}\label{lem:7-2-184bis}\itshape
设 $\{\mathcal F_i\colon C_i\to
\mathfrak C_{i+1}^{(r)}\}$
是有向系统
$\{I_i\colon C_i\to C_{i+1}\}$、
$\{I'_i\colon C'_i\to C'_{i+1}\}$
之间态射 $\{f_i\colon C_i\to C'_i\}$
的交换误差同伦。则有：

\noindent\textup{(7.2.185.3)}
可将 $\{\mathcal F_i\}$ 提升为两个有向系统
之提升之间的滤过 $A_\infty$ 同伦。
具体说，$\mathcal F_i$ 可提升为
滤过 $A_\infty$ 同态，
使 (7.2.186) 成为滤过
$A_\infty$ 同态的严格等式：
\[
\operatorname{Eval}_{r=0}\circ\widetilde{\mathcal F}_i
=I'_i\circ\widetilde f_i,
\qquad
\operatorname{Eval}_{r=1}\circ\widetilde{\mathcal F}_i
=\widetilde f_{i+1}\circ I_i.
\]
其中 $\widetilde f=\{\widetilde f_i\}$、
$\widetilde{\mathcal F}=\{\widetilde{\mathcal F}_i\}$
分别是 $\{f_i\}$ 与 $\{\mathcal F_i\}$ 的提升。
称 $\widetilde{\mathcal F}$ 为该提升的
\emph{交换误差同伦}\liftIndex{I271}。
\end{statement}

事实上，引理 7.2.184bis 已在
第~\ref{subsec:7-2-9} 小节的
引理~\ref{lem:7-2-129} 证明中得到；
另见引理 7.2.208。

\begin{proposition}{7.2.187}\label{prop:7-2-187}
设 $\{I_i\colon C_i\to C_{i+1}\}$、
$\{I'_i\colon C'_i\to C'_{i+1}\}$
为两个有向系统，$\{f_i\}$ 为它们之间的态射，
$\{\mathcal F_i\colon C_i\to
\mathfrak C_{i+1}^{(r)}\}$
为其交换误差同伦。
给定这两个有向系统的 $A_\infty$ 提升，
并设 $f^0=\{f_i^0\}$、$f^1=\{f_i^1\}$
是 $\{f_i\}$ 在这些提升上的两个
$A_\infty$ 提升。
再假定对 $f^0$ 与 $f^1$，
分别给定 $\{\mathcal F_i\}$
的 $A_\infty$ 提升；
并不要求这两种提升相同。
则 $f^0$ 与 $f^1$ 滤过 $A_\infty$ 同伦。
\end{proposition}

命题~\ref{prop:7-2-187} 的证明
留待陈述下文的定理 7.2.212
与推论 7.2.219 后进行。
届时将以更自然、灵活的条件，
替换这里要求 $f^0,f^1$ 共用
$A_{\bar n,\bar k}$ 交换误差同伦
$\{\mathcal F_i\}$ 的条件。

为了准确陈述定理 7.2.212，
须系统使用同伦的同伦。
以下不同模型的 $\operatorname{Eval}$\liftIndexRangeStart{I272}、
$\operatorname{Incl}$\liftIndexRangeStart{I273} 沿用相同记号，
其意义由上下文区分。

\begin{definition}{7.2.188}\label{def:7-2-188}
设 $C_{r_0,s_0}$（$r_0,s_0=0,1$）
是四个滤过 $A_\infty$ 代数，
包含共同的子代数 $C_0$；
即 $C_0$ 的运算 $\mathfrak m_k$
是各 $C_{r_0,s_0}$ 上运算的限制。
假定每个包含映射
$C_0\subset C_{r_0,s_0}$
都是同伦等价，记五元组为
\[
C=(C_0,\{C_{r_0,s_0}\}_{r_0,s_0}).
\]
若滤过 $A_\infty$ 代数
$\widehat{\mathfrak C}$ 满足以下条件，
则称其为 $[0,1]^2\times C$ 的一个
\emph{模型}\liftIndex{I270}。

\noindent\textup{(7.2.189.1)}
分别存在 $[0,1]\times
(C_0,\{C_{r_0,0},C_{r_0,1}\})$
的模型 $\widehat{\mathfrak C}_{r=r_0}$
（$r_0=0,1$），以及
$[0,1]\times
(C_0,\{C_{0,s_0},C_{1,s_0}\})$
的模型 $\widehat{\mathfrak C}_{s=s_0}$
（$s_0=0,1$）；见定义~\ref{def:7-2-174}。
前者的端点映射记为
$\operatorname{Eval}_{s=s_0}$，
后者的端点映射记为
$\operatorname{Eval}_{r=r_0}$。
各边与角点的位置见图~\ref{fig:7-2-10}。

\par\medskip\noindent\begin{minipage}{\textwidth}
\centering
\setlength{\unitlength}{0.8mm}
\begin{picture}(125,72)
\put(42,18){\framebox(43,38){}}
\put(63.5,37){\makebox(0,0){$\widehat{\mathfrak C}$}}
\put(63.5,61){\makebox(0,0){$\widehat{\mathfrak C}_{s=1}$}}
\put(63.5,13){\makebox(0,0){$\widehat{\mathfrak C}_{s=0}$}}
\put(34,37){\makebox(0,0){$\widehat{\mathfrak C}_{r=0}$}}
\put(94,37){\makebox(0,0){$\widehat{\mathfrak C}_{r=1}$}}
\put(38,60){\makebox(0,0){$C_{r=0,s=1}$}}
\put(89,60){\makebox(0,0){$C_{r=1,s=1}$}}
\put(38,13){\makebox(0,0){$C_{r=0,s=0}$}}
\put(89,13){\makebox(0,0){$C_{r=1,s=0}$}}
\put(15,18){\vector(0,1){35}}
\put(10,36){\makebox(0,0){$s$}}
\put(42,4){\vector(1,0){43}}
\put(90,4){\makebox(0,0){$r$}}
\end{picture}
\par{\renewcommand{\thefigure}{7.2.10}\refstepcounter{figure}
\normalfont\bfseries 图~\thefigure\label{fig:7-2-10}}
\end{minipage}\par\medskip

\noindent\textup{(7.2.189.2)}
存在滤过 $A_\infty$ 同态
\[
\operatorname{Eval}_{r=r_0}\colon
\widehat{\mathfrak C}
\to\widehat{\mathfrak C}_{r=r_0},
\qquad
\operatorname{Eval}_{s=s_0}\colon
\widehat{\mathfrak C}
\to\widehat{\mathfrak C}_{s=s_0}.
\]

\noindent\textup{(7.2.189.3)}
对 $k\ne1$，
$(\operatorname{Eval}_{r=r_0})_k$
与 $(\operatorname{Eval}_{s=s_0})_k$ 均为零；
当 $\beta\ne\beta_0=0$ 时，
它们的 $(1,\beta)$ 分量也均为零。

\noindent\textup{(7.2.189.4)}
图式~\ref{diag:7-2-5} 交换。

\par\medskip\noindent\begin{minipage}{\textwidth}
\centering
\[
\begin{array}{ccc}
\widehat{\mathfrak C}&
\xrightarrow{\operatorname{Eval}_{r=r_0}}&
\widehat{\mathfrak C}_{r=r_0}\\[1ex]
\downarrow\scriptstyle\operatorname{Eval}_{s=s_0}&&
\downarrow\scriptstyle\operatorname{Eval}_{s=s_0}\\[1ex]
\widehat{\mathfrak C}_{s=s_0}&
\xrightarrow{\operatorname{Eval}_{r=r_0}}&
C_{r_0,s_0}
\end{array}
\]
\par{\renewcommand{\thefigure}{7.2.5}\refstepcounter{figure}
\normalfont\bfseries 图式~\thefigure\label{diag:7-2-5}}
\end{minipage}\par\medskip

\noindent\textup{(7.2.189.5)}
图式~\ref{diag:7-2-6} 正合。

\par\medskip\noindent\begin{minipage}{\textwidth}
\centering
\[
\begin{gathered}
\widehat{\mathfrak C}
\xrightarrow{A}
\bigoplus_{s_0=0,1}\widehat{\mathfrak C}_{s=s_0}
\oplus
\bigoplus_{r_0=0,1}\widehat{\mathfrak C}_{r=r_0}
\xrightarrow{B}\\
\bigoplus_{\substack{r_1=0,1\\s_1=0,1}}C_{r_1,s_1}
\longrightarrow0.
\end{gathered}
\]
\par{\renewcommand{\thefigure}{7.2.6}\refstepcounter{figure}
\normalfont\bfseries 图式~\thefigure\label{diag:7-2-6}}
\end{minipage}\par\medskip
其中
\[
A=\bigoplus_{s_0=0,1}
(\operatorname{Eval}_{s=s_0})_1
\oplus
\bigoplus_{r_0=0,1}
(\operatorname{Eval}_{r=r_0})_1,
\]
$B$ 按分量定义。
写
$B=\{B^{s=s_0}_{r_1,s_1}\}
\oplus\{B^{r=r_0}_{r_1,s_1}\}$，并令
\[
B^{s=s_0}_{r_1,s_1}
=\delta_{s_0s_1}
(\operatorname{Eval}_{r=r_1})_1,
\qquad
B^{r=r_0}_{r_1,s_1}
=-\delta_{r_0r_1}
(\operatorname{Eval}_{s=s_1})_1,
\]
其中 $\delta_{ij}$ 是克罗内克符号。

\noindent\textup{(7.2.189.6)}
存在严格的滤过 $A_\infty$ 同态
$\operatorname{Incl}\colon C_0\to
\widehat{\mathfrak C}$，使

\noindent\textup{(7.2.189.6.1)}
$\operatorname{Eval}_{s=s_0}\circ\operatorname{Incl}
=\operatorname{Incl}_r$；
两边均将 $C_0$ 映至
$\widehat{\mathfrak C}_{s=s_0}$。

\noindent\textup{(7.2.189.6.2)}
$\operatorname{Eval}_{r=r_0}\circ\operatorname{Incl}
=\operatorname{Incl}_s$；
两边均将 $C_0$ 映至
$\widehat{\mathfrak C}_{r=r_0}$。

\noindent\textup{(7.2.189.7)}
将系数环归约为 $R$ 后，
$\operatorname{Incl}$ 与各
$\operatorname{Eval}$ 均诱导链同伦等价。
\liftIndexRangeEnd{I272}\liftIndexRangeEnd{I273}
\end{definition}

滤过 $A_{n,K}$ 版本和非滤过版本，
都以相同方式定义。
下一个命题证明 $[0,1]^2\times C$ 模型的存在性。

\begin{proposition}{7.2.190}\label{prop:7-2-190}
设 $C_{r_0,s_0}$（$r_0,s_0=0,1$）
是包含共同子代数 $C_0$ 的四个滤过
$A_\infty$ 代数；$C_0$ 的运算
是各角点代数运算的限制，且每个包含映射
$C_0\subset C_{r_0,s_0}$ 都是同伦等价。
记 $C=(C_0,\{C_{r_0,s_0}\}_{r_0,s_0})$。
设 $\widehat{\mathfrak C}_{r=r_0}$
与 $\widehat{\mathfrak C}_{s=s_0}$
分别是 (7.2.189.1) 中的边模型，
并假定所有对象均为 $\mathfrak G$-有隙。
则存在 $\mathfrak G$-有隙的滤过
$A_\infty$ 代数 $\widehat{\mathfrak C}$，
它与上述边模型合起来构成
$[0,1]^2\times C$ 的模型。
\end{proposition}

\par\noindent\textit{命题~\ref{prop:7-2-190} 的证明。}
令
\[
B\colon
\bigoplus_{s_0=0,1}\widehat{\mathfrak C}_{s=s_0}
\oplus
\bigoplus_{r_0=0,1}\widehat{\mathfrak C}_{r=r_0}
\longrightarrow
\bigoplus_{\substack{r_1=0,1\\s_1=0,1}}C_{r_1,s_1}
\]
如图式~\ref{diag:7-2-6} 所定义。
它诱导 $R$-归约 $\overline B$；
例如
\[
\overline{\widehat{\mathfrak C}}_{s=s_0}
[e,e^{-1}]
=
\frac{\widehat{\mathfrak C}_{s=s_0}}
{\Lambda_{0,\mathrm{nov}}^+
\widehat{\mathfrak C}_{s=s_0}}.
\]
令
\begin{equation}
\overline{\widehat{\mathfrak C}}'
=\ker\overline B,
\qquad
\widehat{\mathfrak C}'
=\overline{\widehat{\mathfrak C}}'
\otimes_R\Lambda_{0,\mathrm{nov}}.
\tag{7.2.191}\label{eq:7-2-191}
\end{equation}
$\widehat{\mathfrak C}'$ 是良定义的链复形。
由 (7.2.189.1)，$B=\overline B
\otimes_R\Lambda_{0,\mathrm{nov}}$。
显然可定义 $\operatorname{Eval}_{s=s_0}$、
$\operatorname{Eval}_{r=r_0}$
使 (7.2.189.5) 成立。
此外还有链复形的正合列
\begin{equation}
\begin{gathered}
0\longrightarrow
\overline{\widehat{\mathfrak C}}'
\xrightarrow{\overline A}
\bigoplus_{s_0=0,1}
\overline{\widehat{\mathfrak C}}_{s=s_0}
\oplus
\bigoplus_{r_0=0,1}
\overline{\widehat{\mathfrak C}}_{r=r_0}
\\
\xrightarrow{\overline B}
\bigoplus_{\substack{r_1=0,1\\s_1=0,1}}
\overline C_{r_1,s_1}
\longrightarrow0.
\end{gathered}
\tag{7.2.192}\label{eq:7-2-192}
\end{equation}
\translatornote{原书此式的可见编号误排为“(30.192)”；
依前后 (7.2.191)、(7.2.193) 及后文引用
改为 (7.2.192)。}
将 $\overline A,\overline B$ 换成 $A,B$，
也有类似的正合列。
除 $\widehat{\mathfrak C}'$ 上的滤过
$A_\infty$ 结构外，
(7.2.189.2)、(7.2.189.3) 容易得到。

定义 $\operatorname{Incl}_{1,\beta_0}\colon
\overline C_0\to
\overline{\widehat{\mathfrak C}}'$，
使它与 $\overline A$ 相容，即
\[
\overline A\circ\operatorname{Incl}_{1,\beta_0}
=\bigoplus_{s_0=0,1}
\operatorname{Incl}^{s=s_0}_{1,\beta_0}
\oplus
\bigoplus_{r_0=0,1}
\operatorname{Incl}^{r=r_0}_{1,\beta_0}.
\]
这里及以下，
$\operatorname{Incl}^{s=s_0}\colon C_0\to
\widehat{\mathfrak C}_{s=s_0}$
与
$\operatorname{Incl}^{r=r_0}\colon C_0\to
\widehat{\mathfrak C}_{r=r_0}$
是定义~\ref{def:7-2-174} 中的
滤过 $A_\infty$ 同态。
$\operatorname{Incl}_{1,\beta_0}\colon
\overline C_0\to
\overline{\widehat{\mathfrak C}}'$
不一定是链同伦等价。
但利用同调代数的标准方法，
可找到包含 $\widehat{\mathfrak C}'$
作为链复形的 $\widehat{\mathfrak C}$，
使
$\operatorname{Incl}_{1,\beta_0}\colon
\overline C_0\to
\overline{\widehat{\mathfrak C}}$
成为链同伦等价。
还可取 $\overline{\widehat{\mathfrak C}}$
为以 $\overline{\widehat{\mathfrak C}}'$
作直和项的自由 $R$ 模，即
\[
\overline{\widehat{\mathfrak C}}
=\overline{\widehat{\mathfrak C}}'
\oplus\overline{\widehat{\mathfrak C}}'',
\qquad
\widehat{\mathfrak C}
=\overline{\widehat{\mathfrak C}}
\otimes_R\Lambda_{0,\mathrm{nov}},
\]
这里的分解仅作为 $R$ 模或
$\Lambda_{0,\mathrm{nov}}$ 模成立，
并非链复形的分解。
利用该分解，把 $A$，
亦即各 $\operatorname{Eval}$，
延拓到 $\widehat{\mathfrak C}$，
使 $A$ 成为链映射。
此时 $A$ 不再单射，
但图式~\ref{diag:7-2-6} 仍正合。

现在定义
$\operatorname{Incl}=\{\operatorname{Incl}_{k,\beta}\}_{k,\beta}
\colon C_0\to\widehat{\mathfrak C}$
以及 $\widehat{\mathfrak C}$ 上的
滤过 $A_\infty$ 结构
$\{\mathfrak m_{k,\beta}\}$，
使 (7.2.189) 成立。
按定义~\ref{def:7-2-63} 中
$(n,k)=(\lVert\beta\rVert,k)$ 的次序 $<$
归纳定义 $\operatorname{Incl}_{k,\beta}$
和 $\mathfrak m_{k,\beta}$。
各边模型和 $C_0$ 上的滤过
$A_\infty$ 结构，
分别记为 $\mathfrak m_*^{s=s_0}$、
$\mathfrak m_*^{r=r_0}$、
$\widehat{\mathfrak m}_*^{C_0}$。
假设 $(\lVert\beta'\rVert,k')<(n,k)$ 时，
$\operatorname{Incl}_{k',\beta'}$
与 $\mathfrak m_{k',\beta'}$
已定义。

\begin{lemma}{7.2.193}\label{lem:7-2-193}
设 $(\lVert\beta\rVert,k)=(n,k)$。
则存在映射
\[
\operatorname{Incl}_{k,\beta}\colon
B\overline C_0[1]\longrightarrow
\overline{\widehat{\mathfrak C}},
\qquad
\mathfrak m_{k,\beta}\colon
B_k\operatorname{Im}
(\operatorname{Incl}_{1,\beta_0})
\longrightarrow
\overline{\widehat{\mathfrak C}}
\]
使 $A_\infty$ 关系在
$B_k\operatorname{Im}
(\operatorname{Incl}_{1,\beta_0})$ 上成立，
并使 $\operatorname{Incl}_{k,\beta}$
同此前各 $\operatorname{Incl}_{k',\beta'}$
一起定义滤过 $A_{n,k}$ 同态。
\end{lemma}

引理后一半虽仅在
$B_k\operatorname{Im}
(\operatorname{Incl}_{1,\beta_0})$
上定义 $\mathfrak m_{k,\beta}$，
仍然有意义：滤过 $A_{n,k}$ 同态
在 $B_k\overline C_0$ 上的定义式，
只在该子空间上用到
$\mathfrak m_{k,\beta}$；见 (7.2.196)。

\begin{proof}
由于
$\widehat{\mathfrak C}_{s=s_0}$
与 $\widehat{\mathfrak C}_{r=r_0}$
分别是 (7.2.189.1) 的区间模型，
由定义~\ref{def:7-2-174} 可知
\begin{equation}
\overline{\operatorname{Eval}}_{r=r_0}
\circ\operatorname{Incl}^{s=s_0}_{k,\beta}
=
\overline{\operatorname{Eval}}_{s=s_0}
\circ\operatorname{Incl}^{r=r_0}_{k,\beta}.
\tag{7.2.194}\label{eq:7-2-194}
\end{equation}
由 (7.2.192) 的正合性，
存在 $\operatorname{Incl}_{k,\beta}$ 使
\begin{equation}
\begin{aligned}
\overline{\operatorname{Eval}}_{r=r_0}
\circ\operatorname{Incl}_{k,\beta}
&=\operatorname{Incl}^{r=r_0}_{k,\beta},\\
\overline{\operatorname{Eval}}_{s=s_0}
\circ\operatorname{Incl}_{k,\beta}
&=\operatorname{Incl}^{s=s_0}_{k,\beta}.
\end{aligned}
\tag{7.2.195}\label{eq:7-2-195}
\end{equation}
$\operatorname{Incl}$ 是滤过 $A_{n,k}$ 同态的条件
写成
\begin{equation}
\begin{aligned}
&\bigl(\mathfrak m_{k,\beta}
\circ\operatorname{Incl}_{1,\beta_0}^{\otimes k}\bigr)
(\mathbf x)
+\bigl(\mathfrak m_{1,\beta_0}
\circ\operatorname{Incl}_{k,\beta}\bigr)(\mathbf x)
-\bigl(\operatorname{Incl}_{k,\beta}
\circ\widehat{\mathfrak m}^{C_0}_{1,\beta_0}\bigr)
(\mathbf x)\\
&\quad=
-\sum
\mathfrak m_{\ell,\beta'}
\bigl(\operatorname{Incl}_{k_1,\beta_1}
(\mathbf x_a^{\ell:1}),\ldots,
\operatorname{Incl}_{k_\ell,\beta_\ell}
(\mathbf x_a^{\ell:\ell})\bigr)\\
&\qquad\quad+
\sum\operatorname{Incl}_{k_1,\beta_1}
\circ\widehat{\mathfrak m}^{C_0}_{k_2,\beta_2}
(\mathbf x).
\end{aligned}
\tag{7.2.196}\label{eq:7-2-196}
\end{equation}
右侧第一和取遍
$a,\ell,\beta',k_i,\beta_i$，满足
$k_1+\cdots+k_\ell=k$、
$\beta'+\sum_i\beta_i=\beta$，
并排除
$(\ell,\beta')=(k,\beta)$
及 $(k_i,\beta_i)=(k,\beta)$。
第二和取遍
$k_1+k_2=k+1$、
$\beta_1+\beta_2=\beta$，
并排除 $(k_i,\beta_i)=(k,\beta)$。
这里
\[
\Delta^{\ell-1}(\mathbf x)
=\sum_a
\mathbf x_a^{\ell:1}\otimes\cdots\otimes
\mathbf x_a^{\ell:\ell}.
\]
(7.2.196) 唯一确定
$\mathfrak m_{k,\beta}$ 在
$B_k\operatorname{Im}
(\operatorname{Incl}_{1,\beta_0})$ 上的值。

现证明该子空间上的 $A_{n,k}$ 关系，即
\begin{equation}
\mathfrak m_{k,\beta}\circ
\widehat{\mathfrak m}_{1,\beta_0}
+\mathfrak m_{1,\beta_0}\circ
\mathfrak m_{k,\beta}
=-\sum
\mathfrak m_{k_1,\beta_1}
\circ\widehat{\mathfrak m}_{k_2,\beta_2}.
\tag{7.2.197}\label{eq:7-2-197}
\end{equation}
和取遍 $k_1+k_2=k+1$、
$\beta_1+\beta_2=\beta$，
且 $(k_i,\beta_i)\ne(k,\beta)$。
为证 (7.2.197)，使用 (7.2.196)。
对 $\mathbf x\in B_k\overline C_0[1]$，
得到
\begin{equation}
\begin{aligned}
0=\sum_{\ell',\beta'}
\sum_{\ell'',\beta''}
\sum_{k_i,\beta_i}
\mathfrak m_{\ell',\beta'}
\circ\widehat{\mathfrak m}_{\ell'',\beta''}
\circ\bigl(&\operatorname{Incl}_{k_1,\beta_1}
(\mathbf x_a^{\ell:1}),\ldots,\\
&\operatorname{Incl}_{k_\ell,\beta_\ell}
(\mathbf x_a^{\ell:\ell})\bigr),
\end{aligned}
\tag{7.2.198}\label{eq:7-2-198}
\end{equation}
其中对所有满足
$\ell'+\ell''=\ell+1$、
$\sum_i k_i=k$ 的
$\ell',\beta',\ell'',\beta'',
(k_1,\beta_1),\ldots,(k_\ell,\beta_\ell)$
及 $a$ 求和。
只取其中至少有一项
$(k_i,\beta_i)\ne(1,\beta_0)$ 的部分。
固定这样的一组
$(k_1,\beta_1),\ldots,
(k_\ell,\beta_\ell),a$，
再对满足
$\ell'+\ell''=\ell+1$、
$\beta'+\beta''
=\beta-\sum_i\beta_i$
的 $\ell',\ell'',\beta',\beta''$ 求和；
由归纳假设，此和为零。
故 (7.2.198) 表明，
当所有 $(k_i,\beta_i)=(1,\beta_0)$ 时，
右端相应的和也为零。
这推出 $B_k\operatorname{Im}
(\operatorname{Incl}_{1,\beta_0})$
上的 (7.2.197)。引理得证。
\end{proof}

\begin{lemma}{7.2.199}\label{lem:7-2-199}
存在
\[
\mathfrak m^I_{k,\beta}\in
\operatorname{Hom}
(B_k\overline{\widehat{\mathfrak C}},
\overline{\widehat{\mathfrak C}})
\]
使：

\noindent\textup{(7.2.200.1)}
$\mathfrak m^I_{k,\beta}$ 延拓
引理~\ref{lem:7-2-193}
在 $B_k\operatorname{Im}
(\operatorname{Incl}_{1,\beta_0})$
上定义的 $\mathfrak m_{k,\beta}$。

\noindent\textup{(7.2.200.2)}
$\overline{\operatorname{Eval}}_{s=s_0}
\circ\mathfrak m^I_{k,\beta}
=\mathfrak m^{s=s_0}_{k,\beta}
\circ\overline{\operatorname{Eval}}_{s=s_0}^{\otimes k}$，
其中 $\mathfrak m^{s=s_0}_{k,\beta}$
是 $\widehat{\mathfrak C}_{s=s_0}$
上的滤过 $A_{n,K}$ 结构。

\noindent\textup{(7.2.200.3)}
$\overline{\operatorname{Eval}}_{r=r_0}
\circ\mathfrak m^I_{k,\beta}
=\mathfrak m^{r=r_0}_{k,\beta}
\circ\overline{\operatorname{Eval}}_{r=r_0}^{\otimes k}$，
其中 $\mathfrak m^{r=r_0}_{k,\beta}$
是 $\widehat{\mathfrak C}_{r=r_0}$
上的滤过 $A_{n,K}$ 结构。
\end{lemma}

\begin{proof}
令 $\mathfrak m^{s=s_0,r=r_0}_{k,\beta}$
为角点 $C_{s=s_0,r=r_0}$ 上的滤过
$A_{n,K}$ 结构。由图式~\ref{diag:7-2-5}
的交换性，
\[
\begin{aligned}
\overline{\operatorname{Eval}}_{r=r_0}
\circ\mathfrak m^{s=s_0}_{k,\beta}
\circ\overline{\operatorname{Eval}}_{s=s_0}^{\otimes k}
&=\mathfrak m^{s=s_0,r=r_0}_{k,\beta}
\circ\overline{\operatorname{Eval}}_{r=r_0}^{\otimes k}
\circ\overline{\operatorname{Eval}}_{s=s_0}^{\otimes k}\\
&=\overline{\operatorname{Eval}}_{s=s_0}
\circ\mathfrak m^{r=r_0}_{k,\beta}
\circ\overline{\operatorname{Eval}}_{r=r_0}^{\otimes k}.
\end{aligned}
\]
故引理由图式~\ref{diag:7-2-6}
的正合性和
$\operatorname{Incl}_{1,\beta_0}$
的分裂单射性得出。
\end{proof}

注意这些 $\mathfrak m^I_{k,\beta}$
未必满足 $A_{n,K}$ 关系。
下面将修正它们，并在
$B\overline{\widehat{\mathfrak C}}$
上定义 $\mathfrak m_{k,\beta}$。
令
\[
\mathfrak K=
\bigcap_{r_0=0,1}
\ker\bigl(\overline{\operatorname{Eval}}_{r=r_0}\bigr)
\cap
\bigcap_{s_0=0,1}
\ker\bigl(\overline{\operatorname{Eval}}_{s=s_0}\bigr)
\subset\overline{\widehat{\mathfrak C}}.
\]
它是 $\overline{\widehat{\mathfrak C}}$
的子复形。定义
\[
o_{k,\beta}
(\overline{\widehat{\mathfrak C}})
=\sum_{\substack{k_1+k_2=k\\\beta_1+\beta_2=\beta}}
\mathfrak m^I_{k_1,\beta_1}
\circ\widehat{\mathfrak m}^I_{k_2,\beta_2}
\in\operatorname{Hom}
(B_k\overline{\widehat{\mathfrak C}},
\overline{\widehat{\mathfrak C}}).
\]
则有
\begin{align}
\delta_1\bigl(o_{k,\beta}
(\overline{\widehat{\mathfrak C}})\bigr)
&=0,
\tag{7.2.201.1}\label{eq:7-2-201-1}\\
o_{k,\beta}
(\overline{\widehat{\mathfrak C}})
&\in\operatorname{Hom}
(B_k\overline{\widehat{\mathfrak C}},
\mathfrak K),
\tag{7.2.201.2}\label{eq:7-2-201-2}\\
o_{k,\beta}
(\overline{\widehat{\mathfrak C}})
\circ\operatorname{Incl}_{1,\beta_0}^{\otimes k}
&=0.
\tag{7.2.201.3}\label{eq:7-2-201-3}
\end{align}
这里 $\delta_1$ 是由
$\overline{\mathfrak m}_1$
与 $\overline{\widehat{\mathfrak m}}_1$
诱导的
$\operatorname{Hom}
(B_k\overline{\widehat{\mathfrak C}},
\overline{\widehat{\mathfrak C}})$
上的微分。
(7.2.201.1) 源于归纳假设；
(7.2.201.2) 源于 (7.2.200.2)、
(7.2.200.3)；
(7.2.201.3) 源于 (7.2.200.1)
及 $\mathfrak m_{k,\beta}$
在 $B_k\operatorname{Im}
(\operatorname{Incl}_{1,\beta_0})$
上满足 $A_\infty$ 关系。

由于映射
\[
\operatorname{Hom}
(B_k\overline{\widehat{\mathfrak C}},\mathfrak K)
\longrightarrow
\operatorname{Hom}
(B_k\operatorname{Im}
(\operatorname{Incl}_{1,\beta_0}),\mathfrak K)
\]
在 $\delta_1$ 上同调诱导同构，
由 (7.2.201) 可取
$\operatorname{Cor}\in
\operatorname{Hom}
(B_k\overline{\widehat{\mathfrak C}},\mathfrak K)$，
使它在 $B_k\operatorname{Im}
(\operatorname{Incl}_{1,\beta_0})$ 上为零，
且
\[
o_{k,\beta}
(\overline{\widehat{\mathfrak C}})
+\delta_1\operatorname{Cor}=0.
\]
显然
$\mathfrak m_{k,\beta}
=\mathfrak m^I_{k,\beta}
+\operatorname{Cor}$
具有所需性质。
命题~\ref{prop:7-2-190} 得证。
\hfill$\square$

\begin{remark}{7.2.202}\label{rem:7-2-202}
可以证明
\[
H^d(\mathfrak K)
\cong H^2([0,1]^2;\partial[0,1]^2)
\otimes H^{d-2}
(\overline{\widehat{\mathfrak C}}).
\]
因此
\[
[o_{k,\beta}
(\overline{\widehat{\mathfrak C}})]
\in H^0\!\left(
\operatorname{Hom}
(B_k\overline{\widehat{\mathfrak C}},\mathfrak K)
\right)
\cong H^{-2}\!\left(
\operatorname{Hom}
(B_k\overline{\widehat{\mathfrak C}},
\overline{\widehat{\mathfrak C}})
\right).
\]
这与上述代数障碍论证的几何解释相符。
\end{remark}

取命题~\ref{prop:7-2-190} 的滤过
$A_{n,K}$ 版本，便可得到滤过
$A_{n,K}$ 代数 $C$ 的
$[0,1]^2\times C$ 模型。

\begin{remark}{7.2.203}\label{rem:7-2-203}
也可迭代 $[0,1]\times C$ 模型的构造，
得到 $[0,1]^2\times C$ 的模型。
例如可取
\[
\operatorname{Poly}
([0,1],\operatorname{Poly}([0,1],C))
=\operatorname{Poly}([0,1]^2,C)
\]
（见第 4.6.5 小节），或取
\[
(C^{[0,1]})^{[0,1]}
=C\oplus C[-1]\oplus C[-1]\oplus C[-2]
\]
（比较引理 4.2.25）。
构造细节从略。
这些模型对应于所有角点均有
$C_{r_0,s_0}=C_0=C$ 的情形。
此时自然假定存在映射
\[
\operatorname{Incl}_r\colon
\widehat{\mathfrak C}_{r=r_0}
\longrightarrow\widehat{\mathfrak C},
\qquad
\operatorname{Incl}_s\colon
\widehat{\mathfrak C}_{s=s_0}
\longrightarrow\widehat{\mathfrak C},
\]
\translatornote{原书此处把 $\operatorname{Incl}_r$、
$\operatorname{Incl}_s$ 的定义域交换；
由随后的 (7.2.189bis2) 及图式 7.2.7--7.2.8
按边模型的实际型别写出。}
并满足下述条件。

\noindent\textup{(7.2.189bis1)}
当 $k\ne1$ 时，
$(\operatorname{Incl}_r)_k
=(\operatorname{Incl}_s)_k=0$；
若 $\beta\ne\beta_0$，
其 $(1,\beta)$ 分量也均为零。

\noindent\textup{(7.2.189bis2)}
复合映射
\[
\widehat{\mathfrak C}_{r=r_0}
\xrightarrow{\operatorname{Incl}_r}
\widehat{\mathfrak C}
\xrightarrow{\operatorname{Eval}_{r=r_0}}
\widehat{\mathfrak C}_{r=r_0},
\qquad
\widehat{\mathfrak C}_{s=s_0}
\xrightarrow{\operatorname{Incl}_s}
\widehat{\mathfrak C}
\xrightarrow{\operatorname{Eval}_{s=s_0}}
\widehat{\mathfrak C}_{s=s_0}
\]
均为恒等映射。

\noindent\textup{(7.2.189bis3)}
图式~\ref{diag:7-2-7} 交换。

\par\medskip\noindent\begin{minipage}{\textwidth}
\centering
\[
\begin{array}{ccc}
C&\xrightarrow{\operatorname{Incl}_s}&
\widehat{\mathfrak C}_{r=r_0}\\[1ex]
\downarrow\scriptstyle\operatorname{Incl}_r&&
\downarrow\scriptstyle\operatorname{Incl}_r\\[1ex]
\widehat{\mathfrak C}_{s=s_0}&
\xrightarrow{\operatorname{Incl}_s}&
\widehat{\mathfrak C}
\end{array}
\]
\par{\renewcommand{\thefigure}{7.2.7}\refstepcounter{figure}
\normalfont\bfseries 图式~\thefigure\label{diag:7-2-7}}
\end{minipage}\par\medskip

\noindent\textup{(7.2.189bis4)}
图式~\ref{diag:7-2-8} 交换。

\par\medskip\noindent\begin{minipage}{\textwidth}
\centering
\[
\begin{array}{ccc@{\qquad\qquad}ccc}
\widehat{\mathfrak C}_{r=r_0}&
\xrightarrow{\operatorname{Incl}_r}&
\widehat{\mathfrak C}&
\widehat{\mathfrak C}_{s=s_0}&
\xrightarrow{\operatorname{Incl}_s}&
\widehat{\mathfrak C}\\[1ex]
\downarrow\scriptstyle\operatorname{Eval}_{s=s_0}&&
\downarrow\scriptstyle\operatorname{Eval}_{s=s_0}&
\downarrow\scriptstyle\operatorname{Eval}_{r=r_0}&&
\downarrow\scriptstyle\operatorname{Eval}_{r=r_0}\\[1ex]
C&\xrightarrow{\operatorname{Incl}_r}&
\widehat{\mathfrak C}_{s=s_0}&
C&\xrightarrow{\operatorname{Incl}_s}&
\widehat{\mathfrak C}_{r=r_0}
\end{array}
\]
\par{\renewcommand{\thefigure}{7.2.8}\refstepcounter{figure}
\normalfont\bfseries 图式~\thefigure\label{diag:7-2-8}}
\end{minipage}\par\medskip

由 (7.2.189bis) 取
$\operatorname{Incl}
=\operatorname{Incl}_r\circ\operatorname{Incl}_s$，
便得 (7.2.189.6)。
不过，我们不把 (7.2.189bis)
列为 $[0,1]^2\times C$ 模型定义的一部分，
因为下一小节会用到各
$C_{r_0,s_0}$ 随 $(r_0,s_0)$ 改变的情形。
\end{remark}

为陈述下一个定理，
须给出定义~\ref{def:7-2-182}
的参数化形式。
区间按需要以 $r$ 或 $s$ 参数化；
指定参数时写作 $[0,1]_r$ 或 $[0,1]_s$。
相应的区间模型也分别记作
$\operatorname{Incl}_r$ 与
$\operatorname{Incl}_s$，以资区分。
若 $(C,\{\mathfrak m_{k,\beta}\})$
是滤过 $A_{n,K}$（分别地 $A_\infty$）代数，
就说它是 $A_{n,K}$（分别地 $A_\infty$）
\emph{类}的。
同态、同伦、$[0,1]\times C$ 模型等的
“$A_{n,K}$ 类”“$A_\infty$ 类”
也同样定义。
